% !TEX root = ../jsr_paper.tex
% Section 1: Introduction (fragment; parent declares \section{Introduction})
% Numbers authoritative from _DRAFTING_BRIEF.md. Honest framing: parity not superiority;
% supersonic-validated to M4.33; hypersonic = demonstrated capability, exploratory.

\subsection{Background and Motivation}

Sounding rockets, supersonic missile testbeds, and student-built research
vehicles routinely traverse Mach~1 to Mach~5 during boost and coast. Mission
designers in this class need trajectory tools that predict apogee, peak dynamic
pressure, peak Mach number, and static and dynamic stability across the full
ascent profile with engineering accuracy, and increasingly they need those
tools to be open so that the underlying coefficient models can be audited,
extended, and reproduced. The two most widely used open-source rocket
trajectory simulators, OpenRocket~\cite{niskanen2009} and
RocketPy~\cite{rocketpy2021}, both descend from the Barrowman engineering
aerodynamic method~\cite{barrowman1967} and are reliable below approximately
Mach~1.1. Above the sonic line the underlying constant-coefficient lookup
tables, the incompressible skin-friction transformation, and the absence of
post-shock local-flow corrections cause the predicted drag and stability
derivatives to diverge from measured behavior, frequently producing apogee
errors well in excess of twenty percent on high-supersonic and hypersonic
sounding-rocket missions.

Commercial software fills part of this gap. RASAero~II~\cite{rogers2015}
couples a six-degree-of-freedom integrator to a Missile DATCOM-style coefficient
generator and is widely used in the high-power and university-class rocketry
communities. Its source code is closed and its underlying coefficient model is
not externally auditable by the user. The semi-empirical lineage typified by
Missile DATCOM and the Aeroprediction code family supports missile design at
industrial scale~\cite{moore2001,mooremcinvillehymer2002}, and cross-comparisons
such as that of Sooy and Schmidt against measured data~\cite{sooy2005} document
its maturity; but that lineage is likewise closed and proprietary. No
open-source rocket trajectory simulator presently offers externally benchmarked
supersonic aerodynamic prediction coupled to an end-to-end, reproducible
flight-corpus validation.

The objective of the present work is to close that gap within the OpenRocket
fork OpenRocket-Plus: a single open-source, six-degree-of-freedom trajectory
simulator whose aerodynamic subsystems are each traceable to a published
wind-tunnel, ballistic-range, free-flight, or computational fluid-dynamics
benchmark, and whose integrated apogee predictions are validated against an
externally selected corpus of instrumented flights spanning the supersonic
envelope. We validate end-to-end accuracy to Mach~4.33 against flight data and,
separately and exploratorily, assess how far the same method reaches into the
hypersonic regime on historical vehicles.

\subsection{Prior Work}

Barrowman's master's thesis~\cite{barrowman1967} established the lift and
center-of-pressure framework that anchors essentially every open-source rocket
simulator written since 1967. Niskanen~\cite{niskanen2009} implemented and
extended that framework in OpenRocket, adding component-level drag bookkeeping
and a Mach-dependent base-drag table valid through the transonic peak.
RocketPy~\cite{rocketpy2021} provides a Python-language alternative also rooted
in the Barrowman engineering pipeline, with comparable subsonic fidelity and a
similar transonic ceiling. Both tools treat the freestream Mach number as the
governing aerodynamic state uniformly along the body, an approximation that
degrades as upstream shocks and expansions begin to process the local flow.

Rogers' RASAero~II~\cite{rogers2015} couples a Missile DATCOM-style coefficient
generator to a six-degree-of-freedom trajectory integrator and remains the most
widely used non-OpenRocket prediction tool in high-power and university
rocketry. The closed-source semi-empirical lineage against which the present
work positions itself, Missile DATCOM and the Aeroprediction
series~\cite{moore2001,mooremcinvillehymer2002,sooy2005}, handles post-shock
local conditions internally as part of its coefficient generation but exposes
neither its source nor those intermediate states to the user. Figure~\ref{fig:methods_hierarchy}
places these methods in a common hierarchy: the present work occupies the
open-source engineering tier between the subsonic Barrowman/OpenRocket pipeline
and the closed-source industrial Missile DATCOM/Aeroprediction codes, extending
auditable open-source prediction across the supersonic envelope.

On the open-source-software side, recent precedent in this journal supports the
publication of validated open aerospace software. The arcjetCV
study~\cite{quintart2025} established that open-source aerospace software
accompanied by a documented validation campaign and a permanent code archive is
within the scope of the \textit{Journal of Spacecraft and Rockets}, and a 2025
study of a low-cost amateur roll-control system~\cite{lowcostroll2025}
established that supersonic amateur-rocket hardware work anchored on
open-source trajectory simulation fits within that scope as well. The present
work combines those two precedents: open-source supersonic aerodynamics with a
reproducible flight-corpus validation.

\begin{figure}[t]
  \centering
  \includegraphics[width=\columnwidth]{aerodynamic_methods_hierarchy}
  \caption{Hierarchy of aerodynamic prediction methods for rocket trajectory
  simulation. The present work (OpenRocket-Plus) occupies the open-source
  engineering tier between the subsonic Barrowman/OpenRocket pipeline and the
  closed-source Missile DATCOM/Aeroprediction lineage, extending externally
  benchmarked open-source prediction across the supersonic envelope.}
  \label{fig:methods_hierarchy}
\end{figure}

\subsection{Gap Statement}

The combined prior literature leaves three concrete gaps. First, no open-source
rocket trajectory simulator distributes post-shock local-flow conditions to its
component-level aerodynamic calculators; freestream Mach is applied uniformly
along the body, which is increasingly incorrect above approximately Mach~1.5.
Second, the engineering submodels underlying OpenRocket and RocketPy, including
incompressible skin friction, $\cos^2\Lambda_{LE}$ Ackeret fin wave drag,
constant-coefficient base drag, and the absence of hypersonic blending, remain
largely unchanged from their 1960s and 1970s antecedents and lack independent
benchmarks above the transonic peak. Third, no open-source aerodynamic code has
been subjected to end-to-end integrated flight-corpus validation against an
externally selected set of instrumented supersonic flights, with the residual
decomposed into bias and variance terms and compared on a paired basis against a
commercial baseline.

\subsection{Contributions}

The present work addresses each gap with one corresponding contribution.

\begin{enumerate}
  \item \textbf{A shock-geometry pre-pass architecture.} A once-per-timestep
  nose-to-tail surface march publishes post-shock local conditions, namely local
  Mach number, static pressure, and static temperature, to every component-level
  drag and stability calculator downstream of a single shared seam. The pre-pass
  is inert at subsonic Mach and is verified bit-for-bit against the
  Taylor--Maccoll cone solution and the Prandtl--Meyer function at body
  shoulders. We are explicit about what this contribution does and does not buy:
  its value is \textit{local-flow fidelity} (correct post-shock conditions for
  fin loads and stability) and an architectural seam that enables the downstream
  supersonic submodels, rather than a large change in integrated apogee, which is
  dominated by lower-Mach drag and to which the pre-pass contributes only a
  fraction of a percentage point.

  \item \textbf{An externally benchmarked supersonic subsystem replacement.}
  Each underlying Barrowman-era engineering submodel is replaced and
  independently benchmarked against published wind-tunnel, ballistic-range,
  free-flight, or computational fluid-dynamics data, with the governing
  equation, regime of applicability, and mean absolute percent error reported
  per subsystem. Where a model decision is corpus-calibrated rather than
  literature-anchored, we disclose it as such.

  \item \textbf{A reproducible flight-corpus validation.} An externally selected
  corpus of 25 instrumented flights spanning Mach~0.54 to~4.33 enables
  end-to-end integrated validation of the present method and a direct paired
  comparison against RASAero~II. On this corpus OpenRocket-Plus achieves a mean
  signed apogee error of $-0.38\%$ with all 25 flights within $\pm10\%$ of
  measured altitude; the corresponding RASAero~II predictions, as recorded by
  Rogers, give a mean signed error of $+2.46\%$ with 22 of 25 within $\pm10\%$.
  The paired difference in absolute error is not statistically significant
  (Wilcoxon signed-rank $p=0.615$; $95\%$ bootstrap confidence interval on the
  mean paired difference $[-2.16, +0.96]$ percentage points, straddling zero), so
  the honest claim is \emph{parity} with this version-locked commercial baseline,
  not superiority. The corpus, analysis scripts, and source code are released
  openly under a permanent digital object identifier~\cite{rfd_zenodo}.
\end{enumerate}

Two scope qualifications frame these contributions honestly. The headline
validation corpus was selected externally by Rogers for an independent RASAero~II
comparison rather than curated by the present authors, so its accuracy statistics
are an outcome-independent validation result; the recorded RASAero~II values are
Rogers' own predictions and are not fresh, independently rerun pre-flight
estimates. Separately, although the underlying compressible-flow infrastructure
demonstrates a method capability reaching Mach~7, we do \emph{not} claim
hypersonic flight validation in the headline: high-Mach behavior is assessed
exploratorily on a full, uncurated set of historical sounding rockets, where
motor and geometry reconstruction uncertainty dominates the residual.

\subsection{Paper Organization}

Section~II presents the shock-geometry pre-pass architecture and establishes the
analytical-reference verification that anchors the novelty claim. Section~III
documents the underlying compressible-flow infrastructure: atmosphere,
viscosity, the effective ratio of specific heats, the smoothed compressibility
factor through Mach~1, and the closed-form shock and expansion relations.
Sections~IV and~V catalog the drag and stability submodels respectively, each
with its governing equation, regime of applicability, and published benchmark.
Section~VI rolls up the externally benchmarked subsystems with explicit
verification-and-validation methodology and discloses the corpus-calibrated model
decisions. Section~VII validates against four primary published
computational fluid-dynamics comparators, with a fifth study as an independent
corroborating source. Section~VIII reports the headline result:
end-to-end integrated validation against the 25-flight corpus, with
bias--variance decomposition, per-regime breakdown, paired RASAero~II
comparison, and distribution and normality characterization, together with the
decontaminated prospective holdout that addresses the
partly in-sample nature of two calibrated constants. Section~IX reports the
exploratory high-Mach sounding-rocket assessment over the full historical set,
presented without cherry-picking. Section~X discloses known limitations with
quantitative magnitude and root cause where identified, including a four-flight
operational sensitivity sweep, and Section~XI concludes
with prioritized future work.
